\section{Optimal Power Flow Models}
\label{sec:opf}

\subsection{Problem Form and Runtime Scope}

This chapter covers the OPF family built on top of canonical AC/DC network
projection.  The implemented runtime problem classes are:
\begin{itemize}
  \item nonlinear AC/DC OPF for the native primal-dual and parity-IPM routes,
  \item DC OPF with a true-QP model and LP fallback backends,
  \item reactive-power optimization over discrete taps and shunts by
  outer-approximation / coordinate-search heuristics.
\end{itemize}
The runtime solver family is selected from option flags and DC-network presence.

\subsection{Code Map}

The current implementation is concentrated in:
\begin{itemize}
  \item \texttt{src/optimal_power_flow/ac\_opf.cpp} for top-level AC OPF routing and native reduced-space solve orchestration,
  \item \texttt{src/optimal_power_flow/parity\_formulation.cpp} and \texttt{src/optimal_power_flow/parity\_ipm.cpp} for the parity full-space NLP route,
  \item \texttt{src/optimal_power_flow/dc\_opf.cpp} for DC OPF model build and backend selection,
  \item \texttt{src/optimal_power_flow/reactive\_power\_opt.cpp} for the four-phase outer-approximation / coordinate-search workflow,
  \item \texttt{include/hacdcpf/engine} forwarding headers for LCQP, dual simplex, HiGHS, Gurobi, and related engine facilities supplied by the \texttt{mipsolvers} dependency.
\end{itemize}
There is no \texttt{src/engine} directory in this repository.

\subsection{Pseudocode Map}

The relevant code-aligned workflows are:
\begin{itemize}
  \item Algorithms A8--A10 in Section~\ref{sec:algorithms} for AC OPF routing, parity-NLP construction, and parity-IPM solve,
  \item Algorithm A11 for DC OPF,
  \item Algorithm A12 for reactive-power optimization,
  \item Section~\ref{sec:engine-solver-system} for the forwarding engine surface used by DC OPF and RPO result statistics.
\end{itemize}

\subsection{Current Status}

The AC OPF family has multiple implemented routes, but backend coverage differs
by path.  Parity and native nonlinear routes are the main AC/DC engines; DC OPF
builds a specialized QP/LP model; RPO currently runs an OA/coordinate-search
workflow and reuses branch-and-cut-style statistics fields for counters.

\subsection{Top-Level Solver Routing}

The public \texttt{solve\_ac\_opf} path selects among three solver families.
A key implementation rule is automatic DC routing: if the case has DC buses or
VSC converters and \texttt{enable\_primal\_dual=false}, the options are
internally overridden to
\texttt{enable\_primal\_dual=true, use\_parity\_ipm=true}.

\begin{longtable}{L{0.25\linewidth}L{0.26\linewidth}L{0.37\linewidth}}
\toprule
Path & Trigger & Current behavior \\
\midrule
\endhead
Fast dispatch + PF fallback (AC-only) &
\texttt{enable\_primal\_dual=false} and no forced DC parity routing &
Drops DC and converter subsystems, solves economic dispatch, then runs AC PF projection. \\
Native primal-dual OPF &
\texttt{enable\_primal\_dual=true}, \texttt{use\_parity\_ipm=false} &
Uses a custom sparse KKT loop with AC equations, DC-bus equations, and converter power-balance equations. \\
Parity primal-dual IPM &
\texttt{enable\_primal\_dual=true}, \texttt{use\_parity\_ipm=true}
or auto-enabled by DC presence &
Builds the parity full-space NLP and solves by Mehrotra-type primal-dual IPM. Optional fallback to native path when enabled. \\
\bottomrule
\end{longtable}

\subsection{Fast Dispatch + PF Fallback (AC-Only Path)}

This path is used for fastest feasible dispatch recovery in AC-only operation.
It solves a clipped economic dispatch first, then projects via AC PF.

\paragraph{Economic dispatch stage.}
\begin{align}
\min_{P_g}\; \sum_{g\in\mathcal{G}}
\left(c_{2g}P_g^2 + c_{1g}P_g + c_{0g}\right),
\end{align}
subject to generator bounds and
\(\sum_g P_g = P_D\) (clamped to feasible range).
For quadratic units, the code uses incremental-cost clipping
\(P_g = \mathrm{clip}((\lambda-c_1)/(2c_2),P_g^{\qual{min}},P_g^{\qual{max}})\).

\paragraph{Feasibility projection stage.}
The dispatch is passed to AC Newton PF; resulting bus injections are split back
across generators at each bus and clamped to individual bounds.
If direct projection fails, one reshaped economic-dispatch attempt is executed
before declaring failure.

\subsection{Native Primal-Dual OPF Formulation}

\subsubsection{Decision Variables}

For the native path, variables are
\begin{align}
x_{nat}=
\begin{bmatrix}
\theta_{\mathcal{N}_{ac}\setminus\mathcal{N}_{sl}} &
v_{\mathcal{N}_{PQ}} &
P_g & Q_g &
v^{\qual{dc}} &
P^{\qual{conv,ac}} & Q^{\qual{conv,ac}} & P^{\qual{conv,dc}}
\end{bmatrix}^{\top}.
\end{align}

The implementation includes all DC bus voltages in \(v^{\qual{dc}}\); for
\texttt{DC\_V} buses, tight bounds are built around the setpoint.

\subsubsection{Objective}
\begin{align}
\min_x f(x) =
\sum_{g\in\mathcal{G}}
\left(c_{2g}P_g^2 + c_{1g}P_g + c_{0g}\right).
\end{align}

\subsubsection{Equality Constraints}

\paragraph{AC equations.}
AC rows reuse the PF residual kernel (with converter table disabled in that
kernel call), then subtract explicit converter variables:
\begin{align}
g_{P,i} &= \left(P_i^{\qual{spec,PF}}-P_i^{\qual{calc}}\right) - \sum_{c\in\mathcal{C}_i} P_c^{\qual{ac}} = 0,
&& i\in\mathcal{N}_{ac}\setminus\mathcal{N}_{sl}, \\
g_{Q,i} &= \left(Q_i^{\qual{spec,PF}}-Q_i^{\qual{calc}}\right) - \sum_{c\in\mathcal{C}_i} Q_c^{\qual{ac}} = 0,
&& i\in\mathcal{N}_{PQ}.
\end{align}

\paragraph{DC equations.}
For each DC bus \(k\):
\begin{align}
g_{k}^{\qual{dc}} = P_{d,k}^{\qual{dc}} + V_k^{\qual{dc}}(G_{dc}V^{\qual{dc}})_k
- \sum_{c\in\mathcal{C}_k^{\qual{dc}}} P_c^{\qual{dc}} = 0.
\end{align}

\paragraph{Converter energy balance.}
For converter \(c\):
\begin{align}
g_c^{\qual{conv}} &= P_c^{\qual{ac}} + P_c^{\qual{dc}} + P_c^{\qual{loss}} = 0, \\
P_c^{\qual{loss}} &= a_c + b_c I_c^{\qual{ac}} + (1-\eta_c)(I_c^{\qual{ac}})^2, \\
I_c^{\qual{ac}} &= \frac{\sqrt{(P_c^{\qual{ac}})^2+(Q_c^{\qual{ac}})^2+\epsilon}}{V_c^{\qual{ac}}},
\quad \epsilon=10^{-6}.
\end{align}

\subsubsection{Bounds and Nonlinear Limits}

Variable bounds are enforced for
\(\theta, v, P_g, Q_g, v^{\qual{dc}}, P^{\qual{conv,ac}}, Q^{\qual{conv,ac}}, P^{\qual{conv,dc}}\).

Nonlinear engineering limits are modeled as
\(h(x)\le 0\) and handled by exterior quadratic penalty terms inside the
merit model:
\begin{align}
h_{ij}^{\qual{from}} &= P_{f,ij}^2 + Q_{f,ij}^2 - (S_{ij}^{\qual{max}})^2 \le 0, \\
h_{ij}^{\qual{to}}   &= P_{t,ij}^2 + Q_{t,ij}^2 - (S_{ij}^{\qual{max}})^2 \le 0, \\
h_c^{\qual{conv}}    &= (P_c^{\qual{ac}})^2 + (Q_c^{\qual{ac}})^2 - (S_c^{\qual{max}})^2 \le 0.
\end{align}

\subsubsection{KKT System and Globalization}

At each inner iteration, the solver builds
\begin{align}
\begin{bmatrix}
H + \delta I & J_g^{\top} \\
J_g & -\delta I
\end{bmatrix}
\begin{bmatrix}
\Delta x \\
\Delta\lambda
\end{bmatrix}
=
-\begin{bmatrix}
\nabla_x L \\
g(x)
\end{bmatrix},
\end{align}
with dynamic regularization \(\delta\), line search, trust/restore logic,
and bounded step clipping.

A barrier term is applied to variable bounds; nonlinear branch/converter limits
enter the merit model through adaptive exterior penalties.

\subsection{Parity Full-Space Formulation}

\subsubsection{Decision Variables}

The parity NLP includes AC, DC, converter, shedding, and enhanced DER variables:
\begin{align}
x_{par}=
\begin{bmatrix}
\theta,\ v^{\qual{ac}},\ P_g,\ Q_g,\ v^{\qual{dc}},\ P^{\qual{conv,ac}},\ Q^{\qual{conv,ac}},\ P^{\qual{conv,dc}}, \\
\Delta P_d,\ \Delta Q_d,\ P^{\qual{ren}},\ Q^{\qual{ren}},\ P^{\qual{stor,ac}},\ Q^{\qual{stor,ac}},\ P^{\qual{stor,dc}}, \\
P^{\qual{dcdc}},\ P^{\qual{flex}}
\end{bmatrix}^{\top}.
\end{align}

\subsubsection{Objective}
\begin{align}
\min_x f_{par}(x) =
&\sum_{g\in\mathcal{G}}\left(c_{2g}P_g^2 + c_{1g}P_g + c_{0g}\right)
+ \mathrm{VOLL}_{eff} S_{base}\sum_i (\Delta P_{d,i}+\Delta Q_{d,i}) \\
&+ \sum_{r\in\mathcal{R}^{\qual{curt}}}
\max\left(0, P_{r}^{\qual{rated}} - P_r\right)c_r^{\qual{curt}}.
\end{align}

\subsubsection{Equality Constraints}

\paragraph{AC full-bus balance (scaled).}
For each AC bus \(i\):
\begin{align}
g_{P,i} &=
\Bigl(P_i^{\qual{calc}} - P_i^{\qual{gen}} - P_i^{\qual{fixed}} - P_i^{\qual{ren}} - P_i^{\qual{stor}}
+ P_i^{\qual{load}} - \Delta P_{d,i} + P_i^{\qual{flex}} - P_i^{\qual{conv,ac}}\Bigr)\,\alpha_P = 0, \\
g_{Q,i} &=
\Bigl(Q_i^{\qual{calc}} - Q_i^{\qual{gen}} - Q_i^{\qual{fixed}} - Q_i^{\qual{ren}} - Q_i^{\qual{stor}}
+ Q_i^{\qual{load}} - \Delta Q_{d,i} - Q_i^{\qual{conv,ac}}\Bigr)\,\alpha_Q = 0.
\end{align}
Here \(P_i^{\qual{fixed}},Q_i^{\qual{fixed}}\) aggregate non-variable injections
(static generators, non-curtailable renewables/PV, VPP, microgrid exchange,
mobile storage).

\paragraph{DC balance with DCDC coupling.}
For each DC bus \(k\):
\begin{align}
g_k^{\qual{dc}} = P_{d,k}^{\qual{dc}} + \sum_{m\neq k} g_{km}
\left((V_k^{\qual{dc}})^2 - V_k^{\qual{dc}}V_m^{\qual{dc}}\right)
- P_k^{\qual{conv,dc}} + P_k^{\qual{dcdc,in}} - P_k^{\qual{dcdc,out}} = 0,
\end{align}
where the OPF decision variable is \(P^{\qual{in}}\)
(\emph{input-side referenced}, i.e.\ the power absorbed from \texttt{bus\_in}).
The output power is a deterministic function of \(P^{\qual{in}}\) including
both efficiency and \(I^2R\) conduction loss:
\begin{align}
P^{\qual{out}}=
\begin{cases}
\eta\, P^{\qual{in}} - r_{eq}(P^{\qual{in}})^2/(V^{\qual{dc}})^2,
& P^{\qual{in}}\ge 0, \\
P^{\qual{in}}/\eta - r_{eq}(P^{\qual{in}})^2/(V^{\qual{dc}})^2,
& P^{\qual{in}}<0.
\end{cases}
\end{align}
The parity \texttt{ConstraintIndex} allocates dedicated DCDC balance slots and
counts them in the equality total, but the current equality evaluator leaves
those dedicated rows unpopulated.  The actual DCDC physics is folded directly
into the two DC bus balance equations:
\begin{align}
g_{\texttt{bus\_in}} &\mathrel{+}= P^{\qual{in}}, &
g_{\texttt{bus\_out}} &\mathrel{-}= P^{\qual{out}}.
\end{align}

\paragraph{DCDC Jacobian and Hessian.}
The first-order derivatives of the DCDC coupling terms are:
\begin{align}
\frac{\partial g_{\texttt{bus\_in}}}{\partial P^{\qual{in}}} &= 1, &
\frac{\partial g_{\texttt{bus\_out}}}{\partial P^{\qual{in}}} &=
-\frac{\partial P^{\qual{out}}}{\partial P^{\qual{in}}}, \\
\frac{\partial P^{\qual{out}}}{\partial P^{\qual{in}}} &=
\begin{cases}
\eta - 2r_{eq}P^{\qual{in}}/(V^{\qual{dc}})^2, & P^{\qual{in}}\ge 0, \\
1/\eta - 2r_{eq}P^{\qual{in}}/(V^{\qual{dc}})^2, & P^{\qual{in}}<0.
\end{cases}
\end{align}
When \(r_{eq}>0\), there is also a cross-derivative with respect to
the input-side DC voltage:
\begin{align}
\frac{\partial g_{\texttt{bus\_out}}}{\partial V^{\qual{dc}}_{\texttt{in}}}
&= \frac{2r_{eq}(P^{\qual{in}})^2}{(V^{\qual{dc}}_{\texttt{in}})^3}.
\end{align}
The Hessian contributions from the \(I^2R\) loss term are:
\begin{align}
\frac{\partial^2 g}{\partial (P^{\qual{in}})^2}
&= \frac{2r_{eq}}{(V^{\qual{dc}})^2}, &
\frac{\partial^2 g}{\partial P^{\qual{in}}\partial V^{\qual{dc}}}
&= \frac{-4r_{eq}P^{\qual{in}}}{(V^{\qual{dc}})^3}, &
\frac{\partial^2 g}{\partial (V^{\qual{dc}})^2}
&= \frac{6r_{eq}(P^{\qual{in}})^2}{(V^{\qual{dc}})^4}.
\end{align}

\paragraph{OPF warm-start convention.}
Note that the PF uses output-referenced \(P^{\qual{ref}}\) while the OPF uses
input-referenced \(P^{\qual{in}}\).  The warm-start path converts between these
reference frames when seeding the OPF initial point from a PF result.

\paragraph{Converter balance.}
For each VSC converter:
\begin{align}
P_c^{\qual{ac}}+P_c^{\qual{dc}}+P_c^{\qual{loss}}=0,
\end{align}
with the same loss model as in the native formulation.

\subsubsection{Inequality Constraints}

Nonlinear inequalities:
\begin{align}
P_{f,ij}^2 + Q_{f,ij}^2 - (S_{ij}^{\qual{max}})^2 &\le 0, \\
P_{t,ij}^2 + Q_{t,ij}^2 - (S_{ij}^{\qual{max}})^2 &\le 0, \\
(P_c^{\qual{ac}})^2 + (Q_c^{\qual{ac}})^2 - (S_c^{\qual{max}})^2 &\le 0, \\
\left(g_{km}V_k^{\qual{dc}}(V_k^{\qual{dc}}-V_m^{\qual{dc}})\right)^2 - (P_{km}^{\qual{max}})^2 &\le 0.
\end{align}

Variable bounds are converted to inequality rows:
\begin{align}
x^{\qual{min}}-x \le 0, \qquad x-x^{\qual{max}}\le 0.
\end{align}

\subsubsection{Primal-Dual IPM System}

The parity solver introduces slack \(z>0\) and multipliers \(\mu>0\):
\begin{align}
g(x) &= 0, \\
h(x)+z &= 0, \\
\nabla f(x)+J_g^{\top}\lambda+J_h^{\top}\mu &= 0, \\
z \circ \mu &= \gamma\mathbf{1}.
\end{align}

A Mehrotra predictor-corrector method is used, with:
\begin{enumerate}
  \item affine predictor step,
  \item adaptive centering parameter
  \(\sigma = (\mu_{aff}/(\mu_{cur}+10^{-16}))^3\) after clamping the ratio to \([0,1]\),
  \item corrector step with second-order compensation,
  \item iterative refinement (2 steps) for conditioning of the KKT solve,
  \item fraction-to-boundary step lengths for primal and dual updates.
\end{enumerate}

\paragraph{Implementation note.}
The current \texttt{parity\_ipm.cpp} uses the Mehrotra predictor-corrector
scheme and implements up to two Gondzio extra correctors.
The KKT system is solved using either Eigen \texttt{SparseLU} for sparse or
\texttt{PartialPivLU} for dense factorizations.

Convergence monitors primal infeasibility, dual stationarity, and average
complementarity. A best-iterate restoration is applied before final failure.

\subsection{Component Coverage by Solver Path}

\begin{longtable}{L{0.30\linewidth}L{0.22\linewidth}L{0.36\linewidth}}
\toprule
Component / feature & Native primal-dual & Parity IPM \\
\midrule
\endhead
\texttt{Generator} \((P_g,Q_g)\) & decision variable & decision variable \\
\texttt{VSCConverter} \((P^{\qual{conv,ac}},Q^{\qual{conv,ac}},P^{\qual{conv,dc}})\) & decision variable & decision variable \\
\texttt{DCBus} voltage & decision variable & decision variable \\
AC branch MVA limits & nonlinear penalty term & nonlinear inequality + slack \\
DC branch power limits & not explicit in native path & nonlinear inequality + slack \\
Load shedding \((\Delta P_d,\Delta Q_d)\) & not modeled & decision variable with VOLL penalty \\
Curtailable renewable / controllable PV & not modeled as variable & decision variable \\
AC storage dispatch & not modeled as variable & decision variable \\
DC storage dispatch & not modeled as variable & decision variable; packed into \texttt{ACOPFResult::pstor\_mw} with \texttt{stor\_map.source\_type} identifying AC vs DC storage \\
DCDC dispatch & not modeled as variable & decision variable + DC coupling \\
Flexible load dispatch & not modeled as variable & decision variable \\
StaticGen / fixed DER aggregates & not explicit variable in native objective model & included as fixed injections \\
\bottomrule
\end{longtable}

\subsection{AC OPF Result JSON Contract}

\texttt{opf\_result\_to\_json} writes the scalar solve summary
\texttt{converged}, \texttt{iterations}, \texttt{outer\_iterations},
\texttt{objective}, \texttt{max\_constraint\_violation},
\texttt{max\_stationarity}, \texttt{status}, and \texttt{solver\_path}.  The
\texttt{solver\_path} token is \texttt{NativeAC}, \texttt{ParityIPM}, or
\texttt{Unknown}.

The core vectors are \texttt{vm}, \texttt{va}, \texttt{vdc}, \texttt{pg\_mw},
\texttt{qg\_mvar}, \texttt{dpd\_mw}, \texttt{dqd\_mvar}, \texttt{pac\_mw}, and
\texttt{qac\_mvar}.  Enhanced dispatch is serialized both as structured objects
and as legacy bare arrays:
\begin{itemize}
  \item \texttt{renewable\_dispatch}: \texttt{pren\_mw}, \texttt{qren\_mvar}, and
  a \texttt{map} array of \texttt{original\_index}/\texttt{source\_type} records;
  \item \texttt{storage\_dispatch}: \texttt{pstor\_mw}, \texttt{qstor\_mvar}, and
  a \texttt{map} array; this packed storage array includes AC and DC storage;
  \item \texttt{dcdc\_dispatch}: \texttt{pdcdc\_mw} and a \texttt{map} array;
  \item \texttt{flex\_dispatch}: \texttt{pflex\_mw} and a \texttt{map} array;
  \item legacy top-level \texttt{pren\_mw}, \texttt{qren\_mvar}, \texttt{pstor\_mw},
  \texttt{qstor\_mvar}, \texttt{pdcdc\_mw}, and \texttt{pflex\_mw} arrays.
\end{itemize}
The \texttt{profiling} object contains \texttt{linear\_solver\_backend},
\texttt{analyze\_calls}, \texttt{factorization\_calls},
\texttt{linear\_solve\_calls}, \texttt{total\_iterations},
\texttt{accepted\_steps}, \texttt{rejected\_steps}, and
\texttt{final\_barrier\_mu}.

\subsection{Fallback Behavior}

\begin{enumerate}
  \item If parity IPM is selected and fails, and
  \texttt{allow\_fallback=true}, the solver retries with
  \texttt{use\_parity\_ipm=false}.
  \item If native primal-dual fails and fallback is allowed,
  it attempts the fast dispatch + AC PF path.
  \item Status strings in \texttt{ACOPFResult::status} record which fallback
  path was taken.
\end{enumerate}

\subsection{DC Optimal Power Flow (\texttt{solve\_dc\_opf})}
\label{sec:dc-opf}

The DC OPF model solves a linearized economic dispatch problem using the
DC power flow approximation.  It is independent of the AC OPF paths above
and is invoked through \texttt{hacdcpf::opf::solve\_dc\_opf}.

\subsubsection{Decision Variables}

\begin{align}
x_{dc} =
\begin{bmatrix}
\theta_{\mathcal{N}\setminus\{sl\}} &
P_g &
P_{f} &
\Delta P_d
\end{bmatrix}^{\top},
\end{align}
where \(P_{f,ij}\) are branch flow variables (optional, depending on
formulation mode) and \(\Delta P_{d,i}\) are per-bus load-shedding variables.

\subsubsection{Network Model}

The DC power flow relationship is:
\begin{align}
P_{f,ij} = \frac{\theta_i - \theta_j}{x_{ij}},
\end{align}
with the bus susceptance matrix \(B\) assembled from branch reactances
\(x_{ij}\).  The slack bus angle is fixed at \(\theta_{sl}=0\).

\subsubsection{Objective Function}

The DC OPF builder constructs both an LP model and a QP model.  The QP model
keeps true quadratic generator costs for QP-capable backends; LP backends use
the linearized model.

\paragraph{LP linearized cost.}
\begin{align}
\min \sum_{g\in\mathcal{G}}
c_g^{\qual{lin}} P_g
+ \sum_i \mathrm{VOLL}_{\qual{eff}} \cdot \Delta P_{d,i} \cdot S_{\qual{base}},
\end{align}
where the linearized coefficient is evaluated at the generator midpoint:
\begin{align}
c_g^{\qual{lin}} = c_{1g} + 2c_{2g}\frac{P_g^{\qual{min}}+P_g^{\qual{max}}}{2}.
\end{align}

\paragraph{QP quadratic cost.}
\begin{align}
\min \frac{1}{2}x^{\top}Qx + c^{\top}x,
\end{align}
where \(Q_{gg} = 2c_{2g}S_{\qual{base}}^2\) and \(c_g = c_{1g}S_{\qual{base}}\).
This mode is used by \texttt{NativeQP} and Gurobi.  The HiGHS path currently
calls the LP adapter; QP support there is marked TODO in the source.

\paragraph{VOLL auto-tuning.}
The load-shedding penalty is auto-tuned as \(1.1\times\) the maximum generator
marginal cost to guarantee load serving priority.

\subsubsection{Constraints}

\paragraph{Power balance (per bus).}
\begin{align}
\sum_{g\in\mathcal{G}_i} P_g - P_{d,i} + \Delta P_{d,i}
= \sum_{j\in\mathcal{N}_i} P_{f,ij},
\quad \forall i.
\end{align}

\paragraph{Flow definition.}
\begin{align}
P_{f,ij} - b_{ij}(\theta_i - \theta_j) = 0,
\quad \forall (i,j)\in\mathcal{E}.
\end{align}

\paragraph{Variable bounds.}
\begin{align}
\theta_{sl} &= 0, \qquad
-\pi \le \theta_i \le \pi, \quad i\ne sl, \\
P_g^{\qual{min}} &\le P_g \le P_g^{\qual{max}}, \\
-r_a^{\qual{max}} &\le P_{f,ij} \le r_a^{\qual{max}}, \\
0 &\le \Delta P_{d,i} \le \max(P_{d,i},\, 0.01/S_{\qual{base}}).
\end{align}

\subsubsection{Solver Backend Selection}

\texttt{DCOPFOptions::solver} defaults to \texttt{DCOPFSolverBackend::Auto}.  The
solver always builds the QP form and then dispatches by backend:
\begin{enumerate}
  \item \textbf{NativeQP} (LCQP solver): true quadratic costs via primal-dual IPM.
  \item \textbf{Gurobi}: commercial QP solver.
  \item \textbf{HiGHS}: open-source LP solver (linearized costs).
  \item \textbf{Native Dual Simplex}: via \texttt{solve\_lp\_with\_basis()}.
  \item \textbf{Auto}: tries \texttt{NativeQP}, then Gurobi, then HiGHS, then native dual simplex.
\end{enumerate}

\subsubsection{Locational Marginal Prices}

LMP extraction is controlled by \texttt{DCOPFOptions::compute\_lmp} (default
\texttt{true}).  When enabled, the solver runs a supporting simplex LP after
the primal solve to recover constraint dual variables; the resulting nodal
prices populate \texttt{DCOPFResult::lmp}.  Set
\texttt{compute\_lmp=false} in high-volume loops (Monte Carlo, FMEA) where
price signals are not needed, to skip the dual extraction overhead.

\textbf{Branch shadow prices} (\texttt{branch\_mu\_lower}/\texttt{branch\_mu\_upper})
are populated \emph{only} when branch-flow variables ($P_{ij}$) are included
in the formulation, i.e.\ when \texttt{include\_branch\_limits=true} (the
default).  When branch limits are disabled or the LP is solved without
explicit $P_{ij}$ variables (e.g.\ after a QP solve without a supporting LP),
both vectors are left zero-filled.  Do \textbf{not} interpret zero values as
``uncongested''; check \texttt{DCOPFOptions::include\_branch\_limits} first.

\textbf{Known limitation — congested branch shadow prices:}
The supporting simplex LP used for dual recovery (\texttt{solve\_lp\_with\_basis})
does not handle bounded-variable LP problems.  When \texttt{include\_branch\_limits=true}
and a branch limit is actively binding (congestion), the native simplex cannot
recover correct $P_{ij}$ box duals; \texttt{branch\_mu\_lower/upper} will be
unreliable (typically zero).  LMP nodal prices remain valid.  A bounded-variable
simplex or an interior-point solver would be required to obtain correct branch
shadow prices under congestion.

\subsection{Reactive Power Optimization (\texttt{solve\_rpo})}
\label{sec:rpo}

The reactive power optimization (RPO) module searches over discrete voltage and
reactive-power controls (transformer on-load tap changers and switchable shunt
steps) with continuous AC OPF evaluations.  The current source implements the
OA/coordinate-search/ternary/pairwise refinement workflow below rather than a
generic branch-and-bound MINLP solve loop.

\subsubsection{Problem Structure}

\paragraph{Discrete variables.}
\begin{itemize}
  \item Transformer OLTC tap positions:
    \(\text{tap ratio} = 1 + (\text{tap\_pos} - \text{tap\_neutral}) \times \text{step\_pct}/100\).
  \item Switchable shunt steps:
    \(B_s = \text{step} \times b_{\text{per\_step}}\) (susceptance per step).
\end{itemize}

\paragraph{Continuous variables.}
Bus voltages and generator reactive outputs, solved via AC OPF at each
discrete configuration.

\subsubsection{Objective Functions}

Selectable objectives:
\begin{align}
\text{MinVoltageDeviation:} &\quad
\sum_i w_{\qual{vdev}} (V_i - V_{\qual{tgt}})^2, \\
\text{MinActiveLoss:} &\quad
\sum_g P_g - \sum_d P_d, \\
\text{Combined:} &\quad
w_{\qual{vdev}} \sum_i (V_i - V_{\qual{tgt}})^2
+ \left(\sum_g P_g - \sum_d P_d\right).
\end{align}

\subsubsection{Algorithm: Four-Phase Outer Approximation}

\paragraph{Phase~0: Sensitivity ranking.}
Perturb each discrete variable by \(\pm 1\) from the baseline (two NLP solves
per variable).  Compute \(|\Delta\text{obj}|\) for each variable and order by
descending sensitivity.

\paragraph{Phase~1: Ternary-search sweep.}
For each variable in sensitivity order, exploit approximate unimodality of the
RPO objective using integer ternary search (\(O(\log R)\) evaluations per
variable instead of \(O(R)\)).  At each evaluation, generate an
outer-approximation cutting plane:
\begin{align}
f(y) \ge f_k + g_k^{\top}(y - y_k),
\end{align}
where \(g_k\) is a finite-difference gradient estimate.

\paragraph{Phase~2: Coordinate descent with OA pruning (up to 3 cycles).}
Sweep each variable and evaluate only candidates where the OA lower bound
\(< \text{incumbent}\).  Add OA cuts after each NLP evaluation.  Terminate if
no improvement this cycle.  The OA lower bound is:
\begin{align}
\max_k f_k + g_k^{\top}(y - y_k).
\end{align}

\paragraph{Phase~3: Pairwise neighborhood polishing.}
Check \(\pm 1\) perturbations of every pair of discrete variables.  OA cuts
prune many candidates before expensive NLP evaluation.

\subsubsection{Inner NLP Solver}

At each discrete configuration, the AC OPF solver (IPM by default, fallback
enabled) is called to solve the continuous sub-problem.  Solution caching by
discrete configuration key avoids redundant NLP evaluations.

\subsubsection{Result Structure}

The result includes per-device tap/shunt adjustments, the number of NLP solves,
OA cuts added, nodes explored, and the final status.  Current source status
strings include \texttt{"Optimal (OA)"}, \texttt{"Time limit (OA)"},
\texttt{"Node limit (OA)"}, and \texttt{"No feasible configuration found (OA)"}.
